\section{Discussion: Alignment, Vernacular Competence, and the Standardization of Judgment}
\label{sec:discussion}

The pattern that emerges across the analysis is a specific normative stance that grants abstract entitlement through therapeutic vocabulary without concrete permission, applied regardless of situational clarity. On r/relationship\_advice, therapeutic concepts function as diagnostic tools that authorize prescriptive conclusions: for instance, identifying ``gaslighting'' warrants recommending people leave their relationships. LLMs deploy the same concepts as a uniform register, present everywhere but discriminating nowhere, so that the diagnostic function is severed from its prescriptive consequence.

\subsection{A Governed Ceiling on Directive Judgment}

The most striking finding is structural: on the 153 posts where at least 70\% of community comments recommend leaving, models do not converge with the community; the divergence widens. This is the inverse of what a system calibrated to situational clarity would produce. If LLM caution tracked genuine uncertainty, models would hedge most on ambiguous cases and converge with community responses on clear ones. Instead, the models apply a ceiling on prescriptive confidence regardless of situational clarity, creating an inverse relationship between community certainty and model directiveness. The ceiling does not reflect judgment about particular cases; it reflects a prior about what kind of speech is permissible. Safety alignment is the most plausible account of this prior, though the design cannot rule out contributions from training-data composition and assistant-style instruction tuning.

This has a consequence for understanding safety alignment. The result is not ``bad advice'' but something more precisely characterizable: a
subject who feels validated but is not authorized to act, encouraged to manage emotional responses rather than change a situation. Where the community's therapeutic vocabulary leads to action (``this is 
gaslighting, therefore leave''), LLM therapeutic vocabulary terminates in self-management---what Illouz calls the empowering potential of therapeutic culture \citep{illouz2008saving}, inverted into a form of 
responsibilization that locates structural harm in the victim's capacity for emotional regulation. What alignment yields, more precisely, is \emph{techne} without \emph{phronesis}. 

Vernacular competence in this community, then, is not primarily about emotional attunement or diagnostic accuracy---LLMs perform adequately on both. It is about what Stevanovic and Peräkylä call deontic authority transfer: the willingness to tell someone what to do, and to provide the specific formulation that makes doing it thinkable. The patterns documented here are consistent with this framework: models less often convert diagnostic recognition into explicitly action-guiding language, and when they do, they hedge the prescription in ways the community does not. Lu's politeness research confirms this has evaluative force: generic advice violates ``sociality obligations,'' the expectation that advisers provide tailored, enactable guidance rather than blanket statements \citep{lu2025politeness}. What appears in these outputs is not just suppressed judgment but a recurring substitution: context-specific prescription gives way to a portable advisory style based on caution, validation, and broadly acceptable speech.

The method developed here makes this normative stance empirically characterizable in ways that existing evaluation approaches often do not. Red-teaming tests whether models produce harmful outputs; benchmarks assess helpfulness against generic standards. Neither specifies the structure of alignment's moral economy, the particular configuration of therapeutic uniformity, prescriptive ceilings, and attention-redirection documented here. 

This is not to say that the community's exit-oriented advice is ``unsafe''; it is a different theory of what advice should do. Vicsek et al.\ find that models struggle when cultural sensitivity conflicts with human rights principles in cross-cultural contexts \citep{vicsek2025crosscultural}. The findings reveal a parallel tension operating within Western contexts: safety alignment encodes a particular cultural position (therapeutic neutrality) that conflicts with other defensible positions, including directive intervention in situations the community recognizes as dangerous.

Defamiliarization, if it is to be more than a tool for criticizing one side, must work in both directions. This is worth stating plainly, because the analysis could be misread as endorsing the community's directiveness as the normative standard. It does not. The argument is diagnostic, not prescriptive: the goal is not to recommend that LLMs behave more like r/relationship\_advice, but to characterize what each system assumes and what each makes visible about the other.

To be clear about the paper's scope: the analysis does not claim that r/relationship\_advice gives better advice, that LLM assistants should mimic the subreddit's bluntness, or that under-directiveness is uniformly harmful. What the design establishes is a patterned structural divergence---not a verdict on which advisory form is preferable. Whether the community's concentrated exit-orientation represents good judgment, a product of platform incentives, or something in between is a question this analysis is not equipped to settle and does not attempt to.

The community's taken-for-granted directiveness is revealed as culturally specific; so too is LLM caution. The community's directiveness rests on assumptions that deserve scrutiny: directive exit advice from strangers, based on one-sided accounts, carries real risks. The community cannot know what the poster omitted, cannot follow up to learn whether its advice helped, and cannot verify that exit is materially possible. High consensus may reflect template-matching to familiar scripts as much as situational clarity. LLM caution, whatever its origins, does not presume these conditions are met. A community whose upvote economy rewards bluntness is not obviously a better model for AI advice systems than a professional counselor; indeed, there are strong grounds for preferring something closer to professional standards in many contexts.

But the reverse is also true: the community's capacity to recognize coercive control patterns, isolation tactics, and grooming represents accumulated practical knowledge that therapeutic neutrality cannot access. That this knowledge is culturally specific does not make it wrong; it makes it situated. What the analysis demonstrates is that aligned LLMs do not simply occupy a different normative position from the community---they occupy a structurally consistent one, consistently forgoing deontic authorization in ways that do not vary with situational clarity. The contribution here is characterizing the structure of that pattern: what it assumes, where it operates, and what it renders invisible. The question is not which moral economy is correct, but what it means that the encounter between them is governed by design choices that operate below the level of explicit argumentation.

The repertoire documented in Section~\ref{sec:results} around therapy, communication, reflection, and self-work encodes an implicit theory of what relationship problems are. If problems are located in the advice-seeker's pattern-recognition or self-worth rather than in another person's behavior, then solutions are correspondingly internal: grow, reflect, communicate better. The community's contrasting repertoire locates some problems in the other person (``this is textbook abuse,'' ``he planned things so you can't leave'' etc.) and correspondingly prescribes external action: exit, not self-work. Both are theories of problems, not merely response styles. Alignment assumes problems are communicatively and therapeutically tractable; the community assumes some are structural, requiring action rather than dialogue.

\subsection{Implications for Evaluation}

Current evaluation frameworks treat advice quality as optimizable against a single standard---typically helpfulness as judged by users or other LLMs. Benchmarks assess whether advice is helpful, harmless, and honest, but not whether it grants permission to act. The operationalization of deontic authority transfer developed here suggests this is a measurable property rather than an ineffable quality. Benchmarks that reward only validation will systematically prefer advice that does not provide the authorization situated communities do. 

As LLMs become prevalent in advice-seeking contexts, their systematic divergence from community-specific practices could gradually reshape those practices. The encounter between AI defaults and platform cultures is not a one-time comparison but an ongoing interaction whose dynamics warrant longitudinal study. Evidence from other consequential domains shows that LLM-mediated drafting can alter institutional outcomes (e.g., public comments), suggesting that repeated exposure to model defaults can re-pattern social practices over time \citep{arsenault2026whose}.

This concern connects to what Kleinberg and Raghavan describe as the risk of algorithmic monoculture: when diverse decision-makers are replaced by systems drawing on shared training pipelines, variance contracts even if average quality appears stable \cite{kleinberg2021algorithmic}. The therapeutic flattening documented here provides empirical texture to this dynamic. If advice-seekers increasingly encounter this compressed distribution, moral diversity may erode not through any singularly harmful response but through the cumulative displacement of situated judgment by averaged defaults.

Current safety evaluation focuses on harms of commission: content that is dangerous, misleading, or offensive. The analysis cannot demonstrate that LLM hedging causes harm. But this hedging concentrates on topics where the community's exit orientation is strongest, such as abuse, safety threats, and controlling behavior. This raises questions about whether harms of omission deserve consideration alongside the harms that alignment currently targets.

\subsection{Limitations}

Several limitations qualify the interpretation of these findings and clarify the scope of the claims advanced here.The four models used here are cost-optimized rather than frontier systems; the patterns observed could reflect reduced capacity rather than safety alignment, and replication with frontier models would 
strengthen mechanistic claims. The persona prompting results reported in Section 4.5 weigh against a purely capacity-based account: a model that produces community-style responses when explicitly prompted but not by default is a model whose defaults have been shaped by something other than incapacity. This does not rule out capacity effects, but it shifts the burden toward alignment-related explanations.

The linguistic metrics used here trade semantic precision for scalability. Pairwise validation on 50 human–LLM response pairs supports their directional adequacy for the present analysis: human responses were judged more certain in 72\% of cases, more leave-oriented in 76\%, and LLM responses more therapeutically framed in 96\%. At the same time, these checks do not establish that the lexicons capture the full semantic or pragmatic structure of individual utterances. Inter-coder reliability was uneven, underscoring the difficulty of drawing sharp boundaries in overlapping and highly skewed advice categories. The measures are therefore best understood as coarse instruments for identifying large-scale directional contrasts in advisory style, and not as exhaustive accounts of how judgment, authorization, or therapeutic framing operate in every response.

The analysis focuses on one subreddit representing a specific slice of American online discourse. The posts analyzed reflect multiple selection processes, so the findings characterize \textit{endorsed community advice}, not representative human advice. A related limitation concerns using this community as a baseline at all: r/relationship\_advice's upvote economy actively selects for certainty and directiveness, which means the community's norms around hedging and prescription are themselves distinctive rather than representative of human advice more broadly. The divergence documented here is partly a function of how concentrated these norms are; a community with more diffuse or ambivalent standards would produce smaller gaps and weaker contrast effects. This does not undermine the analysis---the method requires concentrated norms---but it means the findings characterize the LLM's distance from a particular extreme, not from typical human advice. Training data averaging cannot be definitively separated from safety alignment as mechanisms, though the finding that divergence widens where uncertainty is lowest is more consistent with a governed constraint on prescription than with simple averaging.